% =====================================================================
%  PRÉAMBULE COMMUN — Carnet d'entraînement, Fiche 2 (Nomenclature)
%  (inséré physiquement dans chaque document pour qu'il soit autonome)
% =====================================================================
\documentclass[11pt,a4paper]{article}

% ---------- Encodage & langue ----------
\usepackage[utf8]{inputenc}
\usepackage[T1]{fontenc}
\usepackage{lmodern}
\usepackage[french]{babel}

% ---------- Mathématiques & chimie ----------
\usepackage{amsmath,amssymb,mathtools}
\usepackage{icomma}
\usepackage[version=4]{mhchem}
\usepackage{siunitx}
\sisetup{output-decimal-marker={,},per-mode=symbol}

% ---------- Graphismes & mise en page ----------
\usepackage{xcolor}
\usepackage{colortbl}
\usepackage{tikz}
\usepackage[most]{tcolorbox}
\usepackage{enumitem}
\usepackage{array}
\usepackage{booktabs}
\usepackage{multicol}
\usepackage{fancyhdr}
\usepackage{titlesec}
\usepackage[hidelinks]{hyperref}
\usetikzlibrary{arrows.meta,positioning,calc,shapes.geometric,shapes.misc,
  fit,backgrounds,decorations.pathreplacing}
\usepackage[a4paper,margin=2.2cm,headheight=28pt,headsep=10pt]{geometry}

% =====================================================================
%  PALETTE (charte « Chimie » — mauve/prune du cours)
% =====================================================================
\definecolor{primary}{RGB}{146,95,161}
\definecolor{primarydark}{RGB}{92,55,108}
\definecolor{defcol}{RGB}{0,114,178}
\definecolor{thmcol}{RGB}{0,140,120}
\definecolor{formulacol}{RGB}{198,93,9}
\definecolor{remarkcol}{RGB}{120,94,200}
\definecolor{attncol}{RGB}{200,40,55}
\definecolor{excol}{RGB}{0,135,90}
\definecolor{methcol}{RGB}{90,90,100}
\definecolor{metalcol}{RGB}{198,150,40}
\definecolor{nonmetcol}{RGB}{0,140,150}
\definecolor{oxycol}{RGB}{210,55,45}
\definecolor{hydcol}{RGB}{60,120,210}
\definecolor{catcol}{RGB}{200,70,120}
\definecolor{ancol}{RGB}{40,110,190}
\definecolor{saltcol}{RGB}{120,90,160}

% =====================================================================
%  STYLE COMMUN DES BOÎTES
% =====================================================================
\tcbset{
  boxbase/.style={enhanced,breakable,boxrule=0.9pt,arc=2.5pt,
     left=3mm,right=3mm,top=2.3mm,bottom=2.3mm,
     fonttitle=\bfseries,coltitle=white,
     toptitle=0.6mm,bottomtitle=0.6mm},
}

\newtcolorbox{definition}[1][]{%
  boxbase,colback=defcol!5,colframe=defcol,
  title={\textsc{Définition}\ifx\relax#1\relax\relax\else\textmd{ --- \itshape #1}\fi}}
\newtcolorbox{regle}[1][]{%
  boxbase,colback=thmcol!6,colframe=thmcol,
  title={\textsc{Règle}\ifx\relax#1\relax\relax\else\textmd{ --- \itshape #1}\fi}}
\newtcolorbox{formule}[1][]{%
  boxbase,colback=formulacol!6,colframe=formulacol,
  title={\textsc{Méthode-clé}\ifx\relax#1\relax\relax\else\textmd{ --- \itshape #1}\fi}}
\newtcolorbox{remarque}[1][]{%
  boxbase,colback=remarkcol!6,colframe=remarkcol,
  title={\textsc{Remarque}\ifx\relax#1\relax\relax\else\textmd{ : \itshape #1}\fi}}
\newtcolorbox{attention}[1][]{%
  boxbase,colback=attncol!6,colframe=attncol,
  title={\textsc{Attention}\ifx\relax#1\relax\relax\else\textmd{ : \itshape #1}\fi}}
\newtcolorbox{exemple}[1][]{%
  boxbase,colback=excol!5,colframe=excol,
  title={\textsc{Exemple}\ifx\relax#1\relax\relax\else\textmd{ : \itshape #1}\fi}}
\newtcolorbox{methode}[1][]{%
  boxbase,colback=methcol!7,colframe=methcol,sharp corners,
  title={\textsc{Méthode}\ifx\relax#1\relax\relax\else\textmd{ : \itshape #1}\fi}}
\newtcolorbox{astuce}[1][]{%
  boxbase,colback=primary!6,colframe=primary,
  title={\textsc{Astuce concours}\ifx\relax#1\relax\relax\else\textmd{ : \itshape #1}\fi}}

% ---- Exercice (numéroté) ----
\newtcolorbox[auto counter]{exercice}[1][]{%
  boxbase,colback=primary!4,colframe=primarydark,
  title={\textsc{Exercice}~\thetcbcounter\ifx\relax#1\relax\relax\else\textmd{ --- \itshape #1}\fi}}

% ---- Corrigé (numéroté) ----
\newtcolorbox[auto counter]{corrige}[1][]{%
  boxbase,colback=excol!4,colframe=excol,
  title={\textsc{Corrigé}~\thetcbcounter\ifx\relax#1\relax\relax\else\textmd{ --- \itshape #1}\fi}}

% =====================================================================
%  EN-TÊTES ET PIEDS DE PAGE
% =====================================================================
\pagestyle{fancy}
\fancyhf{}
\fancyhead[L]{\small\textcolor{primarydark}{\textbf{Fiche 2} — Nomenclature}}
\fancyhead[R]{\small\textcolor{primarydark}{\textsc{Carnet d'entraînement}}}
\fancyfoot[C]{\small\thepage}
\renewcommand{\headrulewidth}{0.8pt}
\renewcommand{\headrule}{\hbox to\headwidth{\color{primary}\leaders\hrule height \headrulewidth\hfill}}

\titleformat{\section}{\Large\bfseries\color{primarydark}}{\thesection}{0.6em}{}
\titleformat{\subsection}{\large\bfseries\color{primary}}{\thesubsection}{0.6em}{}
\titleformat{\subsubsection}{\normalsize\bfseries\color{primary!80!black}}{}{0em}{}

% ---- Bandeau de titre réutilisable : \bandeau{Thème}{Sous-titre} ----
\newcommand{\bandeau}[2]{%
  \begin{center}
  \begin{tikzpicture}
  \node[rectangle,rounded corners=7pt,fill=primary,text=white,
        minimum width=16.0cm,minimum height=1.7cm,align=center,inner sep=8pt]
    {{\Large\bfseries #1}\\[3pt]{\normalsize #2}};
  \end{tikzpicture}
  \end{center}
  \vspace{0.3cm}}
\begin{document}
\bandeau{Thème 1 — Ions, valences et nombre d'oxydation}{Tuto : la brique de base de toute la nomenclature}

\noindent Avant de \emph{nommer} ou d'\emph{écrire} un composé, il faut savoir lire ses éléments : sont-ils
métalliques ? quelle charge portent-ils ? à quel \textbf{degré d'oxydation} se trouvent-ils ? Ce thème
entraîne cette lecture. Tout le reste du carnet en dépend.

\section*{1. Métaux \emph{vs} non-métaux : le premier réflexe}
\begin{itemize}[leftmargin=1.5em,topsep=2pt,itemsep=1pt]
  \item Un \textbf{métal} \emph{cède} des électrons $\rightarrow$ il forme un \textbf{cation} (charge $\oplus$).
        Exemples : \ce{Na}, \ce{K}, \ce{Ca}, \ce{Mg}, \ce{Al}, \ce{Fe}, \ce{Cu}, \ce{Zn}, \ce{Ag}, \ce{Pb}, \ce{Sn}.
  \item Un \textbf{non-métal} \emph{capte} des électrons $\rightarrow$ il forme un \textbf{anion} (charge $\ominus$)
        ou se lie par covalence. Exemples : \ce{O}, \ce{H}, \ce{N}, \ce{C}, \ce{S}, \ce{P}, \ce{F}, \ce{Cl}, \ce{Br}, \ce{I}.
\end{itemize}

\begin{definition}[valence]
La \textbf{valence} d'un ion est la \emph{valeur absolue de sa charge}. \ce{Na+} et \ce{Cl-} : valence 1 ;
\ce{Ca^2+} et \ce{SO4^2-} : valence 2 ; \ce{Al^3+} et \ce{PO4^3-} : valence 3.
\end{definition}

\begin{astuce}
Un cation a \emph{perdu} des électrons : son nombre d'électrons vaut $Z - (\text{charge})$. Un anion en a
\emph{gagné} : $Z + (\text{|charge|})$. Ex.\ \ce{Ca^2+} : $20-2=18$ électrons ; \ce{O^2-} : $8+2=10$ électrons.
\end{astuce}

\section*{2. Les ions polyatomiques à connaître (par c\oe ur)}
\begin{center}\footnotesize
\renewcommand{\arraystretch}{1.15}
\begin{tabular}{>{\raggedright\arraybackslash}p{5.0cm} >{\raggedright\arraybackslash}p{5.0cm} >{\raggedright\arraybackslash}p{4.2cm}}
\rowcolor{primary!18}
\textbf{Valence 1} & \textbf{Valence 2} & \textbf{Valence 3} \\
\midrule
\textbf{ammonium} \ce{NH4+} \emph{(seul cation !)} & sulfate \ce{SO4^2-} & phosphate \ce{PO4^3-} \\
nitrate \ce{NO3-} / nitrite \ce{NO2-} & sulfite \ce{SO3^2-} & \\
hydroxyde \ce{OH-} & carbonate \ce{CO3^2-} & \\
acétate \ce{CH3COO-} & chromate \ce{CrO4^2-} & \\
permanganate \ce{MnO4-} & dichromate \ce{Cr2O7^2-} & \\
perchlorate \ce{ClO4-} & thiosulfate \ce{S2O3^2-} & \\
\end{tabular}
\end{center}
\noindent\textbf{Piège de signe :} dans cette liste, \ce{NH4+} est le \emph{seul} cation ; tous les autres
sont des anions.

\section*{3. Le nombre d'oxydation (n.o.)}
Le \textbf{nombre d'oxydation} (noté en chiffres romains, $+\mathrm{II}$, $-\mathrm{II}$, \ldots) est la
charge \emph{fictive} qu'aurait un atome si tous les électrons de chaque liaison étaient donnés à l'atome le
plus électronégatif. On l'obtient avec six conventions, appliquées dans l'ordre :

\begin{regle}[les 6 conventions du n.o.]
\begin{enumerate}[label=(\arabic*),leftmargin=2.0em,topsep=2pt,itemsep=0.5pt]
  \item corps simple (\ce{O2}, \ce{Fe}, \ce{Cl2}) : n.o. \textbf{nul} ;
  \item ion monoatomique : n.o. \textbf{= sa charge} (\ce{Cl-}$\to-\mathrm{I}$, \ce{Ca^2+}$\to+\mathrm{II}$) ;
  \item oxygène : $-\mathrm{II}$ (sauf peroxydes $\ce{O2^2-}$ : $-\mathrm{I}$) ;
  \item hydrogène : $+\mathrm{I}$ (sauf hydrures métalliques \ce{H-} : $-\mathrm{I}$) ;
  \item alcalins $+\mathrm{I}$, alcalino-terreux $+\mathrm{II}$, fluor $-\mathrm{I}$ ;
  \item \textbf{somme des n.o. = charge globale} (0 si la molécule est neutre).
\end{enumerate}
\end{regle}

\begin{methode}[calculer le n.o. d'un élément « central »]
On note $x$ le n.o. inconnu, on remplace \ce{O} par $-\mathrm{II}$ et \ce{H} par $+\mathrm{I}$, on multiplie
chaque n.o. par le nombre d'atomes, puis on écrit : \textbf{somme = charge globale}, et l'on résout.
\end{methode}

\begin{exemple}[le n.o. du manganèse dans \ce{KMnO4} (grand classique du concours)]
\ce{KMnO4} est neutre. Le potassium (alcalin) vaut $+\mathrm{I}$, chaque oxygène $-\mathrm{II}$
(il y en a 4). On pose $x$ = n.o. du manganèse :
\[ (+\mathrm{I}) + x + 4\times(-\mathrm{II}) = 0 \;\Longrightarrow\; x - \mathrm{VIII} + \mathrm{I}=0
   \;\Longrightarrow\; \boxed{x = +\mathrm{VII}}. \]
Le manganèse est donc au degré $+\mathrm{VII}$ (et non $+\mathrm{VIII}$ : c'est le piège classique). Comme
$Z_{\ce{Mn}}=25$, l'atome a alors cédé 7 électrons : il en garde $25-7=18$.
\end{exemple}

\begin{attention}
Un \textbf{métal} n'a jamais un n.o. \emph{négatif} à l'état naturel : on écarte d'office « cuivre $-\mathrm{I}$ »
ou « fer $-\mathrm{II}$ ». Le n.o. sert surtout à distinguer les métaux à \emph{valences multiples}
(\ce{Fe^2+}/\ce{Fe^3+}, \ce{Cu+}/\ce{Cu^2+}, \ce{Sn^2+}/\ce{Sn^4+}, \ce{Pb^2+}/\ce{Pb^4+}).
\end{attention}
\end{document}
